\chapter{{\fontsize{14pt}{21pt}\selectfont\MakeUppercase{Проектирование веб-приложения}}}
\label{cha:design}

\section{Общая архитектура системы}

Архитектура веб-приложения приведена в приложении Б на рисунке~\ref{fig:design-architecture}. На схеме выделены две основные части: подготовка корпуса и серверная часть. Такое разделение используется потому, что документация подготавливается заранее, а пользовательский запрос обрабатывается в момент обращения. Ресурсоёмкие операции очистки, разбиения и векторизации вынесены из пользовательского сценария, поэтому серверная часть работает с уже подготовленным индексом~\cite{martin2017,fowler2002,kleppmann2017}.

Подготовка корпуса читает локальные HTML-страницы официальной документации PostgreSQL по версиям 16, 17 и 18, а также дополнительные учебные материалы. Система удаляет служебные элементы страниц, выделяет содержательные блоки, разбивает текст на фрагменты и вычисляет для них векторные представления с помощью модели \texttt{BAAI/bge-m3}. После этого документы, фрагменты, векторы и связанные метаданные сохраняются в PostgreSQL с расширением pgvector. Если выполнять эти действия во время ответа, задержка запроса зависела бы от размера корпуса и от загрузки модели векторизации.

Во время обработки конкретного запроса выполняются только операции, зависящие от вопроса пользователя: валидация входных данных, анализ формулировки, построение векторного представления вопроса, поиск фрагментов, повторное ранжирование, проверка источников и формирование ответа. Серверная часть не обращается к исходным HTML-файлам при каждом запросе. Она работает с нормализованными документами, фрагментами и векторами, сохраненными в базе данных.

PostgreSQL в архитектуре служит точкой, через которую связаны подготовка корпуса и обработка запроса. После переиндексации база содержит поддерживаемые версии, документы, фрагменты, векторные представления и служебные метаданные. Репозиторий поиска получает вектор вопроса, список допустимых корпусов и выбранную версию, а возвращает кандидаты с заголовком документа, путём раздела, текстом, типом корпуса и расстоянием до вопроса. Это упрощает контроль версии, потому что фильтрация выполняется по записи версии в таблице \texttt{versions}, а для официальных документов дополнительно проверяется адрес источника.

Внешняя языковая модель Groq вынесена за пределы локальной части хранения и поиска. В этой архитектуре она отвечает только за языковое оформление результата по подготовленному контексту. Приложение не передаёт модели весь корпус документации и не использует её как самостоятельный источник сведений о PostgreSQL. Перед вызовом Groq API серверная часть уже выбирает фрагменты, проверяет их версию и оценивает достаточность подтверждений.

Поток данных между компонентами остается проверяемым. Пользовательский интерфейс отправляет вопрос, выбранную версию PostgreSQL и режим ответа. Серверная часть преобразует вопрос в структуру с нормализованным текстом, намерением и техническими терминами. После повторного ранжирования сервис генерации получает не произвольный набор страниц, а ограниченный контекст из отобранных фрагментов. После формирования результата система сохраняет историю запроса и использованные источники, а пользователь получает ответ вместе со сведениями о материалах, на которые он опирается.

В реализации этим компонентам соответствуют FastAPI-приложение, PostgreSQL с расширением pgvector, ORM-модели SQLAlchemy, миграции Alembic и схемы данных Pydantic~\cite{fastapi-docs,pydantic-docs,sqlalchemy-docs,alembic-docs,pgvector}. Генерация итогового текста выполняется через Groq API, а векторные представления формируются локальной моделью \texttt{BAAI/bge-m3}.

\section{Архитектура серверной части}

Серверная часть спроектирована как набор слоев с явным разделением ответственности. Маршруты API принимают HTTP-запросы и передают данные в сервисы. Pydantic-схемы определяют контракт входных и выходных структур. Сервисный слой реализует правила обработки запроса. Репозитории инкапсулируют SQL-запросы и сохранение истории. Слой базы данных содержит ORM-модели и подключение к PostgreSQL.

Роли основных компонентов приведены в таблице~\ref{tab:design-server-components}. Компоненты выбраны по фактической структуре проекта, но таблица описывает их на архитектурном уровне, без избыточного перечисления внутренних файлов.

\begingroup
\centering
\small
\setstretch{1}\selectfont
\setlength{\tabcolsep}{4pt}
\begin{longtable}{|P{0.25\textwidth}|P{0.34\textwidth}|P{0.33\textwidth}|}
\caption{Компоненты серверной части}\label{tab:design-server-components}\\
\hline
\tableheadcell{Компонент} & \tableheadcell{Назначение} & \tableheadcell{Связь с обработкой запроса} \\ \hline
\endfirsthead
\multicolumn{3}{l}{\tablepartlabel{окончание таблицы \thetable}}\\
\hline
\endhead
\endfoot
Маршруты API & Прием запросов, вызов сервисов, возврат ответа клиенту. & Запускают сценарии \texttt{/api/ask}, \texttt{/api/versions}, \texttt{/api/history}, диагностики и переиндексации. \\ \hline
Схемы данных & Описание допустимых полей, типов и ограничений. & Фиксируют краткий, подробный и обучающий режимы и структуру источников. \\ \hline
Сервисы & Бизнес-логика обработки вопроса, подготовки корпуса и истории. & Выполняют анализ вопроса, поиск фрагментов, повторное ранжирование, проверку подтверждений и генерацию. \\ \hline
Репозитории & Доступ к таблицам и SQL-выборкам. & Получают кандидаты поиска, сохраняют запросы и связанные источники. \\ \hline
База данных & Хранение версий, документов, фрагментов, векторных представлений и истории. & Обеспечивает версионную фильтрацию, векторный поиск и трассируемость результата. \\ \hline
Пользовательский интерфейс & HTML-страницы, стили и сценарии браузера. & Передаёт вопрос, версию и режим ответа, отображает ответ, источники и историю. \\ \hline
\end{longtable}
\par\endgroup

Проектные маршруты API приведены в таблице~\ref{tab:design-api-routes}. Кроме маршрутов с префиксом \texttt{/api}, приложение предоставляет HTML-страницы \texttt{/} и \texttt{/history}, а также дублирующие маршруты диагностики \texttt{/health} и \texttt{/health/embeddings} без префикса API.

\begingroup
\centering
\small
\setstretch{1}\selectfont
\setlength{\tabcolsep}{2pt}
\begin{longtable}{|P{0.15\textwidth}|P{0.09\textwidth}|P{0.27\textwidth}|P{0.22\textwidth}|P{0.17\textwidth}|}
\caption{Проектный состав маршрутов API}\label{tab:design-api-routes}\\
\hline
\tableheadcell{Маршрут API} & \tableheadcell{Метод} & \tableheadcell{Назначение} & \tableheadcell{Основные входные данные} & \tableheadcell{Результат} \\ \hline
\endfirsthead
\multicolumn{5}{l}{\tablepartlabel{продолжение таблицы \thetable}}\\
\hline
\endhead
\endfoot
\texttt{/api/ask} & POST & Основной пользовательский запрос к приложению. & \texttt{question}, \texttt{pg\_version}, \texttt{answer\_mode}. & Ответ или обучающая структура; список источников. \\ \hline
\texttt{/api/versions} & GET & Получение поддерживаемых версий PostgreSQL. & Нет обязательных параметров. & Список версий с \texttt{docs\_base\_url} и признаком поддержки. \\ \hline
\texttt{/api/history} & GET & Просмотр истории запросов. & Параметр \texttt{limit}. & Записи истории с вопросом, версией, режимом, статусом, задержкой и источниками. \\ \hline
\texttt{/api/health} & GET & Базовая диагностика состояния API. & Нет обязательных параметров. & \texttt{status}, \texttt{timestamp}. \\ \hline
\makecell[l]{\texttt{/api/health/}\\\texttt{embeddings}} & GET & Диагностика модуля векторных представлений. & Нет обязательных параметров. & Поставщик, модель, размерность, размер пакета, окружение и статус. \\ \hline
\makecell[l]{\texttt{/api/admin/}\\\texttt{reindex}} & POST & Запуск переиндексации корпуса. & Список версий, признаки корпусов, ограничение страниц и \texttt{X-Admin-Key}. & Статус запуска и краткая статистика по версиям. \\ \hline
\end{longtable}
\par\endgroup

\section{Сценарий обработки запроса}

Сценарий обработки запроса начинается в пользовательском интерфейсе. Пользователь вводит вопрос, выбирает версию PostgreSQL и один из трёх режимов ответа. Браузер отправляет на \texttt{/api/ask} только поля \texttt{question}, \texttt{pg\_version}, \texttt{answer\_mode}; отдельного флага для подключения дополнительного корпуса нет.

После валидации серверная часть проверяет, поддерживается ли версия. Затем выполняется анализ вопроса: нормализация, определение намерения, выделение технических терминов и формирование расширенного набора терминов для лексической ветки поиска.

Далее система выбирает типы корпусов. Для краткого и подробного режимов выбирается официальный корпус (\texttt{official}); для обучающего режима выбираются официальный и дополнительный корпуса (\texttt{official}, \texttt{supplementary}). После этого вопрос преобразуется в векторное представление, выполняется поиск фрагментов в PostgreSQL, применяется контроль версии, результаты повторно ранжируются, а затем выполняется проверка достаточности подтверждений. Только после этого вызывается генерация ответа или формируется безопасный отказ.

Диаграмма последовательности приведена в приложении Б на рисунке~\ref{fig:design-uml}. Она показывает взаимодействие пользователя, интерфейса, маршрута API, сервисов, базы данных и генеративного адаптера.

\section{Диаграмма процесса обработки запроса}

Процесс обработки пользовательского запроса представлен в приложении Б на BPMN-диаграмме, приведённой на рисунке~\ref{fig:design-bpmn}. Диаграмма показывает ветвление по режиму ответа, проверку источников и разные варианты завершения: успешный ответ, обучающий результат, безопасный отказ или ошибка валидации.

\section{Модель данных и диаграмма сущность-связь}

Модель данных проектируется вокруг трассируемости результата. Ответ должен быть связан не только с текстом вопроса, но и с версией PostgreSQL, режимом ответа и источниками, которые были использованы при генерации. Поэтому база данных содержит таблицы версий, документов, фрагментов, векторных представлений, истории запросов и источников запроса.

Диаграмма сущность-связь приведена в приложении Б на рисунке~\ref{fig:design-er}. Она отражает фактическую схему, реализованную ORM-моделями и миграциями Alembic.

Состав основных таблиц приведён в таблице~\ref{tab:design-data-model}. Все первичные ключи реализованы типом UUID, структура обучающего результата хранится как JSONB, а векторные представления хранятся в отдельной таблице с типом \texttt{vector(1024)}~\cite{postgres-json-types-current,postgres-uuid-type-current,pgvector}.

Отдельная таблица \texttt{versions} используется потому, что версия PostgreSQL участвует почти в каждом запросе. Документ нельзя рассматривать только как HTML-страницу с адресом: один и тот же раздел документации может существовать в разных ветках PostgreSQL и относиться к разным основным версиям. Отдельная сущность версии задаёт точку привязки для документов и истории запросов. Благодаря этому выбранная пользователем версия используется не как текстовая пометка, а как часть модели данных. Такой подход снижает риск смешения материалов PostgreSQL 16, 17 и 18 при поиске и сохранении результата.

Сущности \texttt{documents}, \texttt{chunks} и \texttt{embeddings} разделены намеренно. Документ хранит сведения о странице или учебном материале: заголовок, адрес источника, контрольную сумму, тип корпуса и связь с версией. Фрагмент хранит часть текста, которая участвует в поиске и передаётся в контекст ответа. Векторное представление вынесено в отдельную таблицу, потому что оно является производным от текста фрагмента и зависит от выбранной модели. Если потребуется пересчитать векторные представления другой моделью или другой размерности, изменяется слой поиска, а не исходные документы и не структура истории запросов.

Разделение документа и фрагмента также нужно для сохранения понятного источника. Пользователь видит не весь HTML-документ, а конкретный раздел, который повлиял на ответ. Для этого у фрагмента хранится \texttt{section\_path}. Если бы система сохраняла только документ целиком, источник был бы менее точным: пользователь получил бы ссылку на страницу, но не понимал бы, какая часть страницы использована. Если бы система сохраняла только фрагменты без связи с документом, было бы сложнее восстановить исходный адрес, заголовок и тип корпуса. Связь \texttt{documents}--\texttt{chunks} сохраняет оба уровня: страницу как источник и фрагмент как минимальную единицу поиска.

История запросов отделена от корпуса, потому что она описывает действия пользователя, а не содержимое документации. Корпус может быть переиндексирован, дополнен или очищен, но запись пользовательского обращения должна сохранять вопрос, выбранную версию, режим ответа, статус и задержку обработки. В таблице \texttt{query\_history} также разделены текстовый ответ и обучающая структура. Для краткого и подробного режимов используется поле \texttt{answer\_text}, а для обучающего режима хранится JSON-структура \texttt{tutorial\_json}. Такое решение соответствует контракту API: разные режимы ответа имеют разную форму результата.

Для трассировки ответа до использованных фрагментов используется таблица \texttt{query\_sources}. Она сохраняет позицию источника в выдаче, оценку релевантности, роль источника, заголовок, адрес, путь раздела и тип корпуса. Даже если связанный фрагмент позже будет удален при переиндексации, в записи источника остаются основные сведения, показанные пользователю. Это важно для проверяемости: история запроса хранит не только итоговый текст, но и объясняет, на какие материалы опиралась система при формировании ответа.

Такая модель данных связывает три уровня работы приложения: корпус документации, обработку текущего запроса и последующую проверку результата. Корпус дает документы, фрагменты и векторные представления. Сервис обработки выбирает часть этих фрагментов по версии и режиму ответа. История фиксирует, какие фрагменты попали в результат. Поэтому проверяемость ответа реализована не только в интерфейсе, но и в структуре базы данных.

\begingroup
\centering
\small
\setstretch{1}\selectfont
\setlength{\tabcolsep}{2pt}
\begin{longtable}{|P{0.14\textwidth}|P{0.32\textwidth}|P{0.25\textwidth}|P{0.19\textwidth}|}
\caption{Проектная модель данных}\label{tab:design-data-model}\\
\hline
\tableheadcell{Таблица} & \tableheadcell{Ключевые поля} & \tableheadcell{Назначение} & \tableheadcell{Связи} \\ \hline
\endfirsthead
\multicolumn{4}{l}{\tablepartlabel{продолжение таблицы \thetable}}\\
\hline
\endhead
\endfoot
\texttt{versions} & \texttt{id}, \texttt{major\_version}, \texttt{docs\_base\_url}, \texttt{is\_supported}, \texttt{loaded\_at}. & Хранение поддерживаемых основных версий PostgreSQL. & Связана с \texttt{documents} и \texttt{query\_history}. \\ \hline
\texttt{documents} & \texttt{id}, \texttt{version\_id}, \texttt{title}, \texttt{source\_url}, \texttt{checksum}, \texttt{corpus\_type}, \texttt{audience\_level}. & Хранение страниц официального и дополнительного корпуса. & Связана с \texttt{versions} и \texttt{chunks}. \\ \hline
\texttt{chunks} & \texttt{id}, \texttt{document\_id}, \texttt{chunk\_index}, \texttt{section\_path}, \texttt{chunk\_text}, \texttt{token\_count}, \texttt{content\_type}, \texttt{pedagogical\_role}. & Хранение текстовых фрагментов документации. & Связана с \texttt{documents}, \texttt{embeddings}, \texttt{query\_sources}. \\ \hline
\texttt{embeddings} & \texttt{id}, \texttt{chunk\_id}, \texttt{embedding\_model}, \texttt{embedding\_dimension}, \texttt{embedding}, \texttt{indexed\_at}. & Хранение векторных представлений фрагментов. & Один фрагмент имеет один вектор. \\ \hline
\makecell[l]{\texttt{query\_}\\\texttt{history}} & \makecell[l]{\texttt{id}, \texttt{selected\_version\_id},\\\texttt{user\_question}, \texttt{mode},\\\texttt{answer\_text}, \texttt{tutorial\_json},\\\texttt{status}, \texttt{latency\_ms}.} & Хранение истории пользовательских запросов. & \makecell[l]{Связана с\\\texttt{versions} и\\\texttt{query\_sources}.} \\ \hline
\makecell[l]{\texttt{query\_}\\\texttt{sources}} & \makecell[l]{\texttt{id}, \texttt{query\_id}, \texttt{chunk\_id},\\\texttt{rank\_position},\\\texttt{similarity\_score},\\\texttt{source\_role}, \texttt{source\_url},\\\texttt{corpus\_type}.} & Трассировка ответа до использованных источников. & \makecell[l]{Связана с\\\texttt{query\_history}\\и \texttt{chunks}.} \\ \hline
\end{longtable}
\par\endgroup

\section{Алгоритм подготовки корпуса документации}

Подготовка корпуса выполняется как отдельный процесс переиндексации. Он может запускаться из командной строки или через административный маршрут \texttt{/api/admin/reindex}. Входными параметрами являются список версий, признаки включения официального и дополнительного корпуса, а также необязательное ограничение числа страниц.

Алгоритм подготовки корпуса приведён в листинге~\ref{listing:design-indexing-algo}. Он отражает проектную логику без привязки к конкретным строкам кода.

\begin{listingblock}
\captionof{listing}{Алгоритм подготовки корпуса документации}
\label{listing:design-indexing-algo}
\begin{minted}{text}
для каждой выбранной версии PostgreSQL:
    получить или создать запись версии в базе данных
    если выбран официальный корпус:
        загрузить HTML-страницы из corpus/postgres/html/<version>
        очистить HTML и выделить структурные блоки
        разбить текст на фрагменты
        вычислить векторные представления фрагментов
        заменить ранее сохраненный официальный корпус этой версии
    если выбран дополнительный корпус:
        загрузить документы Markdown из corpus/tutorial
        проверить метаданные и пригодность к индексации
        разбить текст на фрагменты
        вычислить векторные представления фрагментов
        заменить ранее сохраненный дополнительный корпус этой версии
\end{minted}
\end{listingblock}

В реализованном процессе подготовка и вычисление векторных представлений выполняются до удаления старых данных текущего корпуса. После подготовки данные сохраняются пакетами, а изменения фиксируются короткими транзакциями. Такой порядок снижает риск длительной блокировки базы данных во время переиндексации.

\section{Алгоритм разбиения документации на фрагменты}

Разбиение документации на фрагменты выполняется после очистки HTML и выделения блоков. Парсер удаляет служебные элементы страницы, навигацию, заголовки сайта, формы, сценарии и другие шумовые элементы. Затем из содержательной части извлекаются заголовки, абзацы, элементы списков, блоки кода и таблицы. Для каждого блока сохраняется путь раздела.

Разбиение в текущей реализации выполняется по символам: используется ограничение \texttt{CHUNK\_SIZE=1200} и перекрытие \texttt{CHUNK\_OVERLAP=200}. Это нужно указать отдельно, поскольку в модели данных поле называется \texttt{token\_count}, но граница фрагмента задаётся не токенизатором модели, а длиной строки. При создании фрагмента дополнительно вычисляется приблизительное число токенов регулярным выражением.

Правила разбиения приведены в таблице~\ref{tab:design-chunking-rules}. Они направлены на снижение смыслового шума и сохранение связи фрагмента с разделом документации.

\begingroup
\centering
\small
\setstretch{1}\selectfont
\setlength{\tabcolsep}{4pt}
\begin{longtable}{|P{0.30\textwidth}|P{0.58\textwidth}|}
\caption{Правила разбиения документации на фрагменты}\label{tab:design-chunking-rules}\\
\hline
\tableheadcell{Правило} & \tableheadcell{Назначение} \\ \hline
\endfirsthead
\multicolumn{2}{l}{\tablepartlabel{продолжение таблицы \thetable}}\\
\hline
\endhead
\endfoot
Не смешивать разные пути разделов & Фрагмент должен относиться к одному разделу документации, чтобы источник был понятен пользователю. \\ \hline
Не смешивать разные типы содержимого & Абзацы, таблицы и блоки кода не объединяются в один фрагмент при смене типа блока. \\ \hline
Ограничивать длину фрагмента & Фрагмент должен быть достаточно компактным для векторизации и передачи в контекст генерации. \\ \hline
Использовать перекрытие & Часть предыдущего текста переносится в следующий фрагмент при переполнении, чтобы не потерять связность. \\ \hline
Назначать педагогическую роль & Фрагменты получают роль \texttt{overview}, \texttt{prerequisite}, \texttt{step}, \texttt{example} или \texttt{warning}. \\ \hline
\end{longtable}
\par\endgroup

\section{Механизм поиска фрагментов}

Поиск фрагментов строится как комбинирование семантической и лексической веток. Семантическая ветка использует расстояние между векторным представлением вопроса и векторными представлениями фрагментов в pgvector. Лексическая ветка ищет совпадения расширенных терминов в названии документа, пути раздела и тексте фрагмента. Такой подход учитывает как свободные формулировки пользователя, так и точные технические термины.

На первом этапе выбираются кандидаты по версии PostgreSQL и типу корпуса. Затем кандидаты двух веток объединяются по идентификатору фрагмента.

\Needspace{45mm}
Параметр предварительного ранжирования вычисляется по формуле~(\ref{eq:design-pre-rank}):
\begin{equation}
\label{eq:design-pre-rank}
\text{pre\_rank} = 0{,}74 \cdot \text{semantic\_similarity} + 0{,}26 \cdot \text{lexical\_signal},
\end{equation}
\noindent где
\begin{tabular}[t]{@{}l@{ -- }p{0.67\textwidth}@{}}
$\text{pre\_rank}$ & предварительная оценка релевантности фрагмента;\\
$\text{semantic\_similarity}$ & семантическое сходство фрагмента с пользовательским вопросом, вычисленное по косинусному расстоянию;\\
$\text{lexical\_signal}$ & нормированная лексическая оценка совпадений технических терминов.
\end{tabular}

Формула не является итоговой оценкой качества ответа; она используется только для предварительного отбора кандидатов перед повторным ранжированием.

\section{Учёт версии PostgreSQL при поиске}

Контроль версии реализуется двумя слоями. Первый слой -- SQL-фильтрация по \texttt{Version.major\_version}. Она не позволяет выбрать документ, связанный с другой записью версии. Второй слой -- постфильтрация официальных источников по URL. Для официального корпуса допускаются только источники, содержащие ожидаемый фрагмент \texttt{/docs/<pg\_version>/}. Например, для PostgreSQL 16 источник \texttt{/docs/17/sql-select.html} и источник \texttt{/docs/current/sql-select.html} должны быть исключены.

Дополнительный корпус также привязан к версии через таблицу \texttt{documents}, но его URL может иметь схему \texttt{supplementary://} или внешний источник из обработанного корпуса Markdown. Поэтому URL-проверка применяется именно к официальным источникам, а базовая версионная изоляция обеспечивается связью документа с записью версии.

Проектный смысл контроля версии состоит не в гарантии абсолютной правильности ответа, а в снижении риска смешения фрагментов разных веток документации. Оставшиеся риски связаны с полнотой корпуса и качеством найденных фрагментов, поэтому после поиска выполняется повторное ранжирование и проверка подтверждений.

\section{Повторное ранжирование найденных фрагментов}

Повторное ранжирование корректирует порядок кандидатов после первичного поиска. В отличие от предварительной оценки, оно учитывает больше признаков: семантическое сходство, лексическое пересечение, совпадение с заголовком и путём раздела, совпадение технических терминов, лексический сигнал, тип корпуса, педагогическую роль и специальные корректировки для тем логической репликации, конфигурации, совместимости и сравнительных запросов.

Состав признаков приведён в таблице~\ref{tab:design-reranking-signals}. Приоритет официального корпуса заложен явно: для краткого и подробного режимов выбираются только официальные кандидаты, а для обучающего режима дополнительный корпус может занять ограниченную часть контекста.

\begingroup
\centering
\small
\setstretch{1}\selectfont
\setlength{\tabcolsep}{3pt}
\begin{longtable}{|P{0.24\textwidth}|P{0.60\textwidth}|}
\caption{Сигналы повторного ранжирования}\label{tab:design-reranking-signals}\\
\hline
\tableheadcell{Сигнал} & \tableheadcell{Назначение} \\ \hline
\tableheadcell{1} & \tableheadcell{2} \\ \hline
\endfirsthead
\multicolumn{2}{l}{\tablepartlabel{окончание таблицы \thetable}}\\
\hline
\tableheadcell{1} & \tableheadcell{2} \\ \hline
\endhead
\endfoot
\makecell[l]{\texttt{semantic\_}\\\texttt{similarity}} & Оценивает близость вопроса и фрагмента по векторному расстоянию. \\ \hline
\makecell[l]{\texttt{lexical\_}\\\texttt{overlap}} & Учитывает пересечение токенов вопроса и текста фрагмента. \\ \hline
\makecell[l]{\texttt{title\_section\_}\\\texttt{overlap}} & Повышает оценку фрагментов, где запрос совпадает с заголовком или путём раздела. \\ \hline
\makecell[l]{\texttt{technical\_term\_}\\\texttt{overlap}} & Учитывает совпадение с выделенными техническими терминами. \\ \hline
\texttt{corpus\_bonus} & Повышает приоритет официального корпуса и ограничивает влияние дополнительного корпуса. \\ \hline
\texttt{role\_bonus} & Учитывает педагогическую роль фрагмента, особенно в обучающем режиме. \\ \hline
Тематические бонусы и штрафы & Корректируют оценку для логической репликации, параметров конфигурации, совместимости и сравнительных запросов. \\ \hline
\end{longtable}
\par\endgroup

\section{Проверка достаточности подтверждений}

Проверка достаточности подтверждений выполняется после повторного ранжирования. Её задача -- определить, можно ли передавать найденный контекст генеративной модели как основание для уверенного ответа. Логика проверки не использует внешнюю разметку качества, а опирается на признаки найденных фрагментов: наличие официальных источников, итоговую оценку, семантическое сходство, лексическое пересечение, совпадение с заголовком и технические термины.

Для вопросов по логической репликации проверка дополнительно требует тематические якоря: \texttt{logical replication}, \texttt{publication}, \texttt{subscription}, \texttt{wal\_level}, \texttt{max\_replication\_slots}, \texttt{max\_wal\_senders}, \texttt{publisher}, \texttt{subscriber}. Если официального контекста нет или признаки слишком слабые, система возвращает безопасный отказ. Для текстовых режимов это сообщение о недостатке подтверждённых данных, а для обучающего режима -- структура с пустыми шагами и пояснением ограничения.

\section{Обработка ошибок и безопасные сценарии}

Проектирование обработки ошибок опирается на принцип явного и проверяемого отказа. Если система не может корректно выполнить запрос, пользователь должен получить понятную реакцию, а история запроса должна фиксировать статус. Основные ситуации приведены в таблице~\ref{tab:design-errors}.

\newpage
\begingroup
\centering
\small
\setstretch{1}\selectfont
\setlength{\tabcolsep}{3pt}
\begin{longtable}{|P{0.24\textwidth}|P{0.28\textwidth}|P{0.38\textwidth}|}
\caption{Обработка ошибок и безопасные сценарии}\label{tab:design-errors}\\
\hline
\tableheadcell{Ситуация} & \tableheadcell{Причина} & \tableheadcell{Реакция системы} \\ \hline
\endfirsthead
\multicolumn{3}{l}{\tablepartlabel{окончание таблицы \thetable}}\\
\hline
\tableheadcell{Ситуация} & \tableheadcell{Причина} & \tableheadcell{Реакция системы} \\ \hline
\endhead
\endfoot
Неподдерживаемая версия & \texttt{pg\_version} не входит в список \texttt{16,17,18}. & Ошибка валидации или отказ при разрешении версии. \\ \hline
Недопустимый режим ответа & Значение \texttt{answer\_mode} не соответствует одному из трёх допустимых режимов. & Ошибка валидации Pydantic. \\ \hline
Слишком короткий вопрос & Длина очищенного вопроса меньше трёх символов. & Ошибка валидации. \\ \hline
Нет релевантных фрагментов & Поиск не вернул кандидатов выбранной версии. & Ошибка доменного уровня с кодом 404. \\ \hline
Нет официальных источников & После ранжирования остались только дополнительные источники. & Отказ с кодом 422. \\ \hline
Недостаточно подтверждений & Признаки релевантности официального контекста слишком слабые. & Безопасный отказ без уверенной генерации. \\ \hline
Недоступен Groq API & Нет ключа, неверная модель, ошибка сети или ответ сервиса. & Ошибка доменного уровня с кодом 503 и сохранение неуспешной записи истории. \\ \hline
Ошибка переиндексации & Сбой загрузки, векторизации или сохранения. & Откат транзакции текущего пакета и диагностическое сообщение. \\ \hline
\end{longtable}
\par\endgroup

\section{Выводы}

Во второй главе спроектирована архитектура веб-приложения, включающая предварительную подготовку корпуса и обработку запроса в момент обращения. Описаны серверная часть, маршруты API, диаграмма последовательности, диаграмма процесса, модель данных, алгоритмы подготовки корпуса, разбиения документации на фрагменты, поиска, учёта версии, повторного ранжирования и проверки достаточности подтверждений. Проектные решения согласованы с реализацией и не предполагают функций, отсутствующих в программной части.
